\section[Isoparametric Hypercones]{Two-Generator Algebras and Isoparametric Hypercones}
\label{sec:low-rank}

When the radius and one defining polynomial generate a Laplacian
algebra, isoparametric geometry identifies the minimal level and
reduces the subcalibration inequality to one angular variable.

\subsection{Geometric background: parallel hypersurfaces and cones}
\label{sec:isoparametric-geometric-background}

A connected complete embedded hypersurface of the unit sphere is
\emph{isoparametric} if its principal curvatures are constant.
M\"unzner's theorem \cite{Muenzer80,Muenzer81} gives
\(g\in\{1,2,3,4,6\}\), with alternating multiplicities \(m_1,m_2\)
(equal for odd \(g\)). A degree-\(g\) Cartan--M\"unzner polynomial
describes its parallel family, with regular spherical levels in
\((-1,1)\) and two focal submanifolds at the endpoints. See
Cecil--Ryan \cite[Sections~3.2 and~3.4--3.6]{CecilRyan15} for the
geometry and Chi \cite[Section~4]{Chi20} for the classification.
The cone over a spherical hypersurface has curvatures zero radially
and \(r^{-1}\) times the spherical curvatures otherwise; it is minimal
exactly when its link is. We identify the unique minimal regular level below.

Lawlor's full curvature criterion is necessary and sufficient for area
minimization of minimal isoparametric hypercones
\cite[Section~4.3]{Lawlor91}. Their nonflat strictly minimizing members
also provide geometric models for non-affine entire solutions of the
minimal surface equation, as shown by Simon
\cite[Theorem~6]{Simon89}.


\subsection{The two-generator characterization}

For \(\cA_R(F)=\R[R,F]\), fixing the radius leaves one spherical
invariant. The Mendes--Radeschi description gives:

\begin{theorem}[One-dimensional quotients and the minimal level]
\label{thm:rank-one-classification}
Let \(N\geq3\), and let \(F\) be a nonconstant homogeneous polynomial
on \(\R^N\) with \(\cA_R(F)=\R[R,F]\). Suppose
\[
 \Sigma=\{F=0\}\cap\Sph^{N-1}\neq\varnothing,\qquad
 \nabla F\neq0\quad\hbox{on }\Sigma,
\]
and suppose \(\Sigma\) is minimal in \(\Sph^{N-1}\).
Then \(\{F=0\}\) is the cone over the minimal regular member of a
Cartan--M\"unzner isoparametric family.
In particular, \(\deg F\in\{1,2,3,4,6\}\).
\end{theorem}

\begin{proof}
Put \(e=\deg F\). The case \(e=1\) is a hyperplane, so assume
\(e\geq2\). By Mendes--Radeschi
\cite[Proposition~37(c) and its proof]{MendesRadeschi}, there are
constants \(a\neq0\) and \(b\), with \(b=0\) for odd \(e\), such
that \(\Psi=aF+br^e\) satisfies the normalized Cartan--M\"unzner
equations. Their proof gives a Cartan--M\"unzner polynomial
\(\Phi_g\), of degree \(g\in\{1,2,3,4,6\}\), such that either
\(\Psi=\Phi_g\), or
\begin{equation}\label{eq:rank-one-folded-alternative}
 \Psi=\epsilon(2\Phi_g^2-r^{2g}),\qquad \epsilon\in\{-1,1\}.
\end{equation}
On the unit sphere, \(\Sigma=\{\Psi=b\}\). Its regularity and
\(|\nabla_{\Sph^{N-1}}\Psi|^2=e^2(1-\Psi^2)\) imply \(|b|<1\).

Write \(f_g=\Phi_g|_{\Sph^{N-1}}\) and
\(\Delta\Phi_g=c_g r^{g-2}\). The spherical Cartan--M\"unzner
equations give the mean curvature of its regular levels:
\begin{equation}\label{eq:rank-one-level-mean-curvature}
 \diver_{\Sph^{N-1}}
 \frac{\nabla_{\Sph^{N-1}}f_g}{|\nabla_{\Sph^{N-1}}f_g|}
 =\frac{c_g-g(N-2)f_g}{g\sqrt{1-f_g^2}}.
\end{equation}
Since \(N\geq3\), only one regular level can be minimal.
In the folded case \eqref{eq:rank-one-folded-alternative},
\[
 \Sigma=\{f_g=t\}\cup\{f_g=-t\},\qquad
 t=\sqrt{(1+\epsilon b)/2}\in(0,1).
\]
Both levels are nonempty and regular, so they cannot both be minimal.
Thus \(\Psi=\Phi_g\), and \(\Sigma\) is its minimal regular level
\(f_g=c_g/[g(N-2)]\). The degree restriction follows.
\end{proof}

\subsection{One-variable reduction and classical classification}
\label{sec:isoparametric-reduction}

\subsubsection{Cartan--M\"unzner reduction}

Let \(\Phi\) be a Cartan--M\"unzner polynomial of degree
\(g\in\{1,2,3,4,6\}\) on \(\R^N\).
For \(g\geq2\), write the two
multiplicities as \(m_1,m_2\), reverse the sign of \(\Phi\) if
necessary so that \(1\leq m_1\leq m_2\), and put
\begin{equation}\label{eq:iso-parameters}
 s=m_1+m_2,\qquad t=\frac{m_2-m_1}{s},\qquad
 d=N-2=\frac g2s.
\end{equation}
Here \(d\) is the dimension of the minimal link.  The Euclidean
Cartan--M\"unzner equations are
\begin{equation}\label{eq:CM-euclidean}
 |\nabla\Phi|^2=g^2r^{2g-2},
 \qquad
 \Delta\Phi=\frac{g^2}{2}(m_2-m_1)r^{g-2}.
\end{equation}
Let
\begin{equation}\label{eq:iso-F-z}
 f=\frac{\Phi}{r^g},\qquad z=f-t,
 \qquad F=r^gz=\Phi-tr^g.
\end{equation}
The cone \(C=\{F=0\}\) is the cone over the unique minimal regular
member of the isoparametric family.  It is algebraic: for even \(g\),
\(r^g\) is polynomial, while for odd \(g\) the multiplicities are equal
and \(t=0\).
For \(g\geq2\), the cone is nonflat and smooth away from the origin,
so \(\Sing C=\{0\}\).

The coordinate \(z\) has a direct interpretation in the parallel
isoparametric family.  On its regular part put
\[
 \theta=\frac1g\arccos f\in(0,\pi/g),\qquad f=\cos(g\theta).
\]
Here \(\theta\) is spherical normal distance from the focal submanifold
\(M_+=\{f=1\}\) \cite[Section~4, pp.~14--16]{Chi20}.
The Cartan--M\"unzner equations give
\[
 |\nabla_{\Sph}\theta|=1,\qquad
 \Delta_{\Sph}\theta
 =d\,\frac{\cos(g\theta)-t}{\sin(g\theta)}.
\]
For the normal \(\nabla_{\Sph}\theta\), the last expression is the
divergence, or unaveraged scalar mean curvature.  Thus the minimal leaf
is characterized by
\[
 \cos(g\theta_*)=t,\qquad
 \tan^2(g\theta_*/2)=\frac{m_1}{m_2},\qquad
 z=\cos(g\theta)-\cos(g\theta_*).
\]
In particular \(z\) is a centered isoparametric coordinate, not a
signed distance.  For the local signed distance \(q=\theta-\theta_*\),
\[
 z=-g\sin(g\theta_*)q+O(q^2).
\]
We retain \(z\) because it preserves the polynomial form of the
subcalibration calculation.

On \(\Sph^{N-1}\), equations \eqref{eq:CM-euclidean} give the one-variable
diffusion closure
\begin{align}
 A(z):=|\nabla_{\Sph}z|^2
 &=g^2\bigl(1-(z+t)^2\bigr),                 \label{eq:A-z}\\
 B(z):=\Delta_{\Sph}z
 &=-g^2t-g(g+d)z.                            \label{eq:B-z}
\end{align}
The realizable interval is
\begin{equation}\label{eq:z-interval}
 I_t=[-1-t,1-t].
\end{equation}

\begin{proposition}[One-variable subcalibration equation]
\label{prop:iso-master}
Let \(\ell>0\), let \(h\) be smooth on a neighborhood of \(I_t\), and set
\[
 G=r^\ell h(z),\qquad h(0)=0,\qquad h'(0)>0.
\]
Define
\begin{equation}\label{eq:E-h}
 E(z)=\ell^2h(z)^2+A(z)h'(z)^2
\end{equation}
and
\begin{equation}\label{eq:N-h}
 \mathfrak N_\ell[h]
 =\bigl(Ah''+Bh'+(d+1)\ell h\bigr)E
                  -\frac12Ah'E'.
\end{equation}
Then \(|\nabla G|^2=r^{2\ell-2}E\) and
\(\cL(G)=r^{3\ell-4}\mathfrak N_\ell[h]\), so the sign criterion becomes
\begin{equation}\label{eq:iso-potential-sign}
 G\cL(G)=r^{4\ell-4}h(z)\mathfrak N_\ell[h](z)\geq0.
\end{equation}
If \(E>0\) throughout \(I_t\) and
\begin{equation}\label{eq:h-sign-criterion}
 \frac{h(z)}z>0\quad(z\neq0),\qquad
 z\mathfrak N_\ell[h](z)\geq0\quad(z\in I_t),
\end{equation}
then \(G\) generates a two-sided subcalibration and \(C\) is area minimizing.
Its divergence is
\begin{equation}\label{eq:iso-divergence}
 r\diver X_G=\frac{\mathfrak N_\ell[h]}{E^{3/2}}.
\end{equation}
\end{proposition}

\begin{proof}
In \cref{prop:spherical-quotient-potential}, take \(\varphi=h\).
The spherical gradient
pairing is \(\Gamma_{\Sph}(u,v)=Au'v'\), the spherical Laplacian is
\(Au''+Bu'\), and \(N-1=d+1\). Substitution gives the derivative identities and \eqref{eq:iso-potential-sign}.
For the angular profile \(h\) of \(G\), the multiplier \(a_G\) is
\[
 G=a_GF,\qquad a_G=r^{\ell-g}\frac{h(z)}z,
 \qquad \left.\frac{h(z)}z\right|_{z=0}=h'(0).
\]
The quotient is smooth and positive, while \(E>0\) excludes
critical points off the vertex. Since \(G\) is continuous at the
vertex and has the signs of the polynomial \(F\),
\cref{rem:polynomial-finite-perimeter,cor:gradient-subcalibration}
give area minimality. Formula \eqref{eq:div-normalized-gradient}
gives \eqref{eq:iso-divergence}.
\end{proof}

\subsubsection{Stability and the classical classification}

The squared norm of the second fundamental form of a minimal
isoparametric hypersurface is constant:
\begin{equation}\label{eq:iso-A2}
 |A_\Sigma|^2=d(g-1).
\end{equation}
The constant function attains the first eigenvalue \(-d(g-1)\), so
\eqref{eq:exact-Hardy-gap} gives the Hardy constant
\begin{equation}\label{eq:D-gd}
 D_{g,d}=\frac{(d-1)^2}{4}-d(g-1).
\end{equation}

Wang's stability calculation uses the discriminant
\((N-1)^2-4g(N-2)=4D_{g,d}\)
\cite[Lemma~3.4, p.~216, and proof of Theorem~B, p.~238]{Wang94}.
The following is its sharp Hardy form.

\begin{proposition}[Classical stability in Hardy form]
\label{prop:iso-stability}
For \(g>1\), stability and strict stability of the minimal isoparametric
hypercone are equivalent, and hold precisely when
\begin{equation}\label{eq:iso-stability-threshold}
 d\geq4g-2,\qquad\text{equivalently}\qquad N\geq4g.
\end{equation}
The optimal constant in \eqref{eq:strict-stability-definition} is
\(D_{g,d}\) from \eqref{eq:D-gd}.  In the stable range,
\begin{equation}\label{eq:iso-optimal-gap}
 D_{g,d}=\frac14+
       \frac{d(d-4g+2)}4\geq\frac14.
\end{equation}
If \eqref{eq:iso-stability-threshold} fails, the cone is unstable and hence
not area minimizing.
\end{proposition}

\begin{proof}
The roots of \(4D_{g,d}\), as a quadratic in \(d\), are
\(2g-1\pm2\sqrt{g(g-1)}\).  The lower root lies in \((0,1)\), and the upper
root lies strictly between \(4g-3\) and \(4g-2\).  Thus the integer link
dimensions in question give exactly \eqref{eq:iso-stability-threshold}.
Equation \eqref{eq:iso-optimal-gap} gives the gap; negative second
variation excludes minimization.
\end{proof}

For equal multiplicities, \cref{cor:harmonic-radial-existence}
with \(e=g\), \(\kappa=g(g-1)\) gives this same sharp constant.

\begin{theorem}[{Wang \cite[Theorem~A, p.~238]{Wang94}}]
\label{thm:classical-iso-classification}
For \(g\geq2\), a minimal isoparametric hypercone is area minimizing if
and only if
\begin{equation}\label{eq:classical-iso-list}
 N\geq4g,
 \qquad
 (g,m_1,m_2)\notin\{(2,1,5),(4,1,6)\},
\end{equation}
up to interchanging the multiplicities.  The \(g=1\) member is a
hyperplane and is area minimizing.
\end{theorem}

The nonflat minimizing members are strictly minimizing by Lawlor's criterion
\cite[Section~4.3]{Lawlor91}; see the argument of Tang--Zhang
\cite[Section~2]{TangZhang20}.
The two exceptions are strictly stable with optimal gap \(1/4\),
but not minimizing \cite[Theorem~B, p.~238]{Wang94}.

\subsubsection{The harmonic part of the isoparametric classification}
\label{sec:harmonic-isoparametric-scope}

For harmonic isoparametric polynomials, the dimension threshold in
\cref{ex:harmonic-isoparametric-radial} selects the following
parameters from the classical list.
The Cartan--M\"unzner equations \eqref{eq:CM-euclidean} and the
dimension relation \eqref{eq:iso-parameters} give
\[
 \Delta F=0\ \Longleftrightarrow\ m_1=m_2=m,
 \qquad N=gm+2.
\]
For \(g=4\), the possible equal multiplicities follow from the
homogeneous list of Takagi--Takahashi \cite{TakagiTakahashi72},
Stolz's multiplicity theorem \cite{Stolz99}, and the FKM construction
\cite[Theorem~4.1]{FKM81}. The sextic multiplicity restriction is due
to Abresch \cite{Abresch83}; see also
\cite[Chapter~3]{CecilRyan15}.
The first two columns give the multiplicity data; the last applies
the radial criterion:
\[
\begin{array}{c|c|c}
 g&\text{possible }N&\text{cases satisfying }N\ge4g\\ \hline
 2&2m+2,\ m\ge1&N\ge8\\
 3&5,8,14,26&14,26\\
 4&6,10&\text{none}\\
 6&8,14&\text{none}
\end{array}
\]
The FKM restriction follows directly from the Clifford relations.
A symmetric Clifford system
\(P_0,\ldots,P_a\) on \(\R^{2l}\) has FKM multiplicities
\((a,l-a-1)\), so equality requires \(l=2a+1\).
If \(a\ge2\), the operator \(P_1P_2\) preserves
\(E_+=\ker(P_0-I)\) and squares to \(-I\). Since
\(\dim E_+=l\), this forces \(l\) to be even, a contradiction.
Thus the only equal FKM pair is \((1,1)\), realized by
\(a=1,l=3\); the homogeneous list contributes the additional
quartic pair \((2,2)\). The quartic and sextic rows lie below the stability threshold in
\cref{prop:iso-stability}.


\subsection{Explicit polynomial subcalibrations}
\label{sec:isoparametric-certificates}

The one-variable calculation recovers every minimizing case, including
unequal multiplicities and the required exceptional multipliers.

\begin{theorem}[Explicit polynomial subcalibrations]
\label{thm:iso-polynomial-table}
Orient \(F\) as in \eqref{eq:iso-F-z}, put \(R=r^2\), and assume the
isoparametric multiplicities exist.  Every positive case in
\cref{thm:classical-iso-classification} admits the normalized-gradient
subcalibration generated by the polynomial potential in the following table.
\begin{center}
\renewcommand{\arraystretch}{1.18}
\begin{tabular}{@{}ccl@{}}
\toprule
\(g\)&multiplicities&potential \(G\)\\
\midrule
1&hyperplane&\(F\)\\
2&\(s\geq7\), or \((m_1,m_2)=(3,3)\)&\(RF\)\\
2&\((m_1,m_2)=(2,4)\)&\(\bigl(R+\frac3{10}F\bigr)F\)\\
3&\(m_1=m_2=4\) or \(8\)&\(RF\)\\
4&\(s\geq9\), or \(s=8\) and \(m_1\geq2\)&\(R^2F\)\\
4&\((m_1,m_2)=(3,4)\)&\(\bigl(R^2+\frac7{54}F\bigr)F\)\\
4&\((m_1,m_2)=(2,5)\)&\(\bigl(R^2+\frac7{15}F\bigr)F\)\\
4&\((m_1,m_2)=(1,7)\)&\(R\bigl(R^2+\frac12F\bigr)F\)\\
\bottomrule
\end{tabular}
\end{center}
In every row the total degree of \(G\) is at most ten, independently
of \(N\). Both sextic multiplicity pairs lie below the stability
threshold, as recorded in \cref{sec:harmonic-isoparametric-scope}.
For every nonflat row, the field generated by the displayed polynomial potential proves strict
area minimality, with the mass-gap constant in
\eqref{eq:strict-minimality-constant}.
\end{theorem}
\subsubsection{Potentials obtained by a radial factor}

We first take \(h(z)=z\).  Formula \eqref{eq:N-h} factors completely.

\begin{proposition}[Radial polynomial defect]
\label{prop:radial-Q}
For \(h=z\),
\begin{equation}\label{eq:N-z-Q}
 \mathfrak N_\ell[z]=zQ_\ell(z),
\end{equation}
where
\begin{align}
 Q_\ell(z)={}&g^2(1-t^2)\mu_\ell
 +g^2t(\ell^2-2\nu_\ell)z
 +(\ell^2-g^2)\nu_\ell z^2,              \label{eq:Q-ell}\\
 \mu_\ell={}&(d+1)\ell-\ell^2-gd,\qquad
 \nu_\ell=(d+1)\ell-gd.                  \label{eq:mu-nu}
\end{align}
Moreover,
\begin{equation}\label{eq:mu-window}
 \mu_\ell=D_{g,d}-\left(\ell-\frac{d+1}{2}\right)^2.
\end{equation}
Thus the constant term is nonnegative precisely on the admissible
degree interval obtained from the Jacobi equation.
\end{proposition}

\begin{proof}
Insert \(h=z\), \(h'=1\), \(h''=0\), and
\(E=\ell^2z^2+A\) into \eqref{eq:N-h}; collect powers of \(z\).
Identity \eqref{eq:mu-window} follows from \eqref{eq:D-gd}.
\end{proof}

\paragraph{The cases \texorpdfstring{\(g=2\) and \(g=4\)}{g=2 and g=4}.}

For \(g=2\), \(d=s\), \(\ell=4\), one has
\begin{equation}\label{eq:P2}
 Q_4=8P_2,\qquad
 P_2=(s-6)(1-t^2)-2(s-2)tz+3(s+2)z^2.
\end{equation}
The vertex of this convex quadratic is
\[
 z_*=\frac{(s-2)t}{3(s+2)},
\]
and at the right endpoint \(z_+=2m_1/s\),
\begin{equation}\label{eq:P2-right}
 P_2(z_+)=\frac{16m_1((s+2)m_1-s)}{s^2}>0.
\end{equation}
If \(z_*\leq z_+\), equivalently
\[
 m_1\geq m_{1,*}:=\frac{s(s-2)}{8(s+1)},
\]
the minimum is nonnegative precisely when
\begin{equation}\label{eq:K2}
 K_2:=3(s+2)(s-6)(1-t^2)-(s-2)^2t^2\geq0.
\end{equation}
For fixed \(s\), this quantity increases with \(m_1\).  At the entry value,
\begin{equation}\label{eq:K2-entry}
 K_2(m_{1,*})=
 \frac{3(s-2)(s+2)(s^2-8s-12)}{4(s+1)^2}\geq0
 \quad(s\geq10).
\end{equation}
For \(s=7,8,9\), direct evaluation at \(m_1=1\) gives, respectively,
\[
 \frac{23}{49},\qquad6,\qquad\frac{767}{81}.
\]
If \(z_*>z_+\), the minimum on \(I_t\) is the positive endpoint value
\eqref{eq:P2-right}.  This proves \(Q_4\geq0\) for every \(s\geq7\).
For \(s=6\), the same formula handles \((m_1,m_2)=(3,3)\), while
\((2,4)\) requires the correction in the next subsection.

For \(g=4\), \(d=2s\), \(\ell=8\),
\begin{equation}\label{eq:P4}
 Q_8=128P_4,\qquad
 P_4=(s-7)(1-t^2)-2(s-3)tz+3(s+1)z^2.
\end{equation}
Now
\[
 z_*=\frac{(s-3)t}{3(s+1)},
 \qquad
 P_4(z_+)=\frac{16m_1((s+1)m_1-s)}{s^2}>0.
\]
The vertex enters the interval at \(m_{1,*}=(s-3)/8\).  When it does, the
minimum is controlled by
\begin{equation}\label{eq:K4}
 K_4=3(s+1)(s-7)(1-t^2)-(s-3)^2t^2,
\end{equation}
and
\begin{equation}\label{eq:K4-entry}
 K_4(m_{1,*})=
 \frac{3(s-3)(s+1)(s^2-10s-3)}{4s^2}\geq0
 \quad(s\geq11).
\end{equation}
For \(s=9,10\), the values at \(m_1=1\) are \(52/27\) and \(107/25\).
For \(s=8\) and \(m_1\geq2\),
\(K_4\geq27-52(1/2)^2=14\).  This proves the \(g=4\) entries with a radial factor in
\cref{thm:iso-polynomial-table}.

\paragraph{The case \texorpdfstring{\(g=3\)}{g=3}.}

Here \(m_1=m_2=m\), \(t=0\), \(d=3m\), and we choose \(\ell=5\).  Formula
\eqref{eq:Q-ell} becomes
\begin{equation}\label{eq:g3-Q}
 Q_5(z)=9(6m-20)+16(6m+5)z^2.
\end{equation}
It is positive for \(m=4,8\), the two stable Cartan cases by
\cref{prop:iso-stability} and the classification in
\cite[Section~3.1.1]{NadirashviliTkachevVladut14}.  Since
\(r^{5}z=RF\), this gives the stated polynomial potential.

\subsubsection{Borderline polynomial corrections}

At the endpoint, cancellation of the radial defect's leading odd
term fixes the correction coefficient.

At the critical dimension \(N=4g\), one has \(d=4g-2\).  The natural
endpoint degree is \(\ell=2g\), for which \(\mu_\ell=0\).  Consider the
first transverse correction
\begin{equation}\label{eq:critical-h}
 h=z(1+\lambda z).
\end{equation}

\begin{proposition}[Forced critical correction]
\label{prop:forced-critical-lambda}
For \(t>0\), a two-sided sign near \(z=0\) forces
\begin{equation}\label{eq:lambda-star}
 \lambda=\lambda_*=
 \frac{2gt}{(2g+1)(1-t^2)}.
\end{equation}
For this correction to preserve the sign of \(z\) on all of \(I_t\), it is
necessary and sufficient that
\begin{equation}\label{eq:t-threshold}
 t<\frac{2g+1}{4g+1}.
\end{equation}
\end{proposition}

\begin{proof}
At the critical degree, \(d=4g-2\) and \(\ell=2g\).
Substitution in \eqref{eq:N-h} gives
\[
 z\mathfrak N_{2g}[z(1+\lambda z)]
 =2g^3\bigl[(2g+1)(1-t^2)\lambda-2gt\bigr]z^3+O(z^4).
\]
The lower-order terms vanish. Nonnegativity for both signs of \(z\)
therefore forces the cubic coefficient to vanish, which gives
\eqref{eq:lambda-star}. Since
\(\lambda_*>0\), the only possible sign loss for \(h/z=1+\lambda_*z\)
is at \(z=-1-t\).  The inequality
\(1-\lambda_*(1+t)>0\) is exactly \eqref{eq:t-threshold}.
\end{proof}

For \(g=2,s=6\), condition \eqref{eq:t-threshold} accepts
\((3,3)\) and \((2,4)\), but rejects \((1,5)\).  For \(g=4,s=7\), it
accepts \((3,4)\) and \((2,5)\), but rejects \((1,6)\).
The quadratic correction \(h=z+\lambda z^2\) therefore fails
precisely for the two stable non-minimizing members of the classical
list.

Substituting the three admissible unequal multiplicities in
\eqref{eq:lambda-star} gives the rational coefficients below.
Their global signs are then checked on the exact interval $I_t$.

\begin{proposition}[Three quadratic corrections]
\label{prop:quadratic-corrections}
The following identities hold on the indicated realizable intervals.
\begin{align}
g&=2,\quad (m_1,m_2)=(2,4),\quad
t=\tfrac13,\quad I_t=[-\tfrac43,\tfrac23],\notag\\
h&=z(1+\tfrac3{10}z),\quad \ell=4,\notag\\
z\mathfrak N_4[h]
&=\frac{648}{125}z^4(11z+30);                    \label{eq:24-certificate}\\[1mm]
g&=4,\quad (m_1,m_2)=(3,4),\quad
t=\tfrac17,\quad I_t=[-\tfrac87,\tfrac67],\notag\\
h&=z(1+\tfrac7{54}z),\quad \ell=8,\notag\\
z\mathfrak N_8[h]
&=-\frac{64}{19683}z^4
 (2548z^2-132237z-910926);                       \label{eq:34-certificate}\\[1mm]
g&=4,\quad (m_1,m_2)=(2,5),\quad
t=\tfrac37,\quad I_t=[-\tfrac{10}7,\tfrac47],\notag\\
h&=z(1+\tfrac7{15}z),\quad \ell=8,\notag\\
z\mathfrak N_8[h]
&=-\frac{128}{1125}z^4
 (588z^2-8015z-16550).                            \label{eq:25-certificate}
\end{align}
In every case \(h/z>0\) and \(z\mathfrak N_\ell[h]\geq0\) on \(I_t\).
\end{proposition}

\begin{proof}
The three lower bounds for \(h/z\) are \(3/5\), \(23/27\), and
\(1/3\).  In \eqref{eq:24-certificate}, \(11z+30>0\) on \(I_t\).
The quadratics in parentheses in \eqref{eq:34-certificate} and
\eqref{eq:25-certificate} are convex, and their endpoint values are both
negative; hence they are negative on the entire interval.  The identities
follow by substitution into \eqref{eq:N-h}.
\end{proof}

\begin{proposition}[The \((g,m_1,m_2)=(4,1,7)\) correction]
\label{prop:17-correction}
For \(t=3/4\), \(I_t=[-7/4,1/4]\), \(\ell=10\), and
\(h=z(1+z/2)\),
\begin{equation}\label{eq:17-certificate}
 z\mathfrak N_{10}[h]=z^2P(z),
\end{equation}
where
\begin{equation}\label{eq:P-17}
 P(z)=189z^5+1778z^4+5492z^3+4848z^2-763z+42.
\end{equation}
The polynomial \(P\) is positive on \([-7/4,1/4]\).

\end{proposition}

\begin{proof}
Substitution into \eqref{eq:N-h} gives \eqref{eq:17-certificate}.
The rational Bernstein certificate in \cref{ex:Bernstein-17} proves
\(P>0\) on the whole interval.  Also \(h/z\geq1/8\), and the
potential is \(r^{10}h=R(R^2+F/2)F\).
\end{proof}

\begin{proof}[Proof of \cref{thm:iso-polynomial-table}]
The rows with a radial factor follow from \cref{prop:radial-Q} and the calculations in
the preceding subsection.  The remaining rows follow from
\cref{prop:quadratic-corrections,prop:17-correction}.  In each case \(E>0\) on \(I_t\): if \(z\neq0\), then
\(h\neq0\) and \(E\geq\ell^2h^2>0\); at \(z=0\),
\(E(0)=A(0)h'(0)^2>0\).  Apply \cref{prop:iso-master}.
For every nonflat row the displayed numerator is nonzero, and both phases
are nonempty.  The normalized gradient is smooth away from the
vertex.  Hence \cref{cor:polynomial-strict-minimality} proves strict
minimality directly from the displayed polynomial potential.
\end{proof}

For every nonflat row except the quadratic case in ambient dimension
\(N=9\), the displayed degree satisfies \(\ell\ne(N-1)/2\).
Thus \cref{thm:ground-state-gap} already gives strict stability from
the subcalibration sign. In the exceptional case \(\ell=4=(N-1)/2\),
\cref{prop:iso-stability} gives the optimal Hardy constant \(2\).
The following boundary calculations distinguish the degree gap from
positivity of the Jacobi trace.

\subsubsection{Degenerate traces and quantitative strictness}

\begin{corollary}[Four polynomial potentials with vanishing Jacobi trace]
\label{cor:four-degenerate-certificates}
For the four polynomial potentials in \cref{thm:iso-polynomial-table}
corresponding to
\begin{equation}\label{eq:four-degenerate-list}
 (g,m_1,m_2)=(2,3,3),\ (2,2,4),\ (4,3,4),\ (4,2,5),
\end{equation}
one has \(n=d+1=4g-1\), \(\ell=2g\), and
\begin{equation}\label{eq:degenerate-trace-gap}
 \beta_G\equiv0,\qquad J_C(|\nabla G|^{-1})=0,\qquad
 \inf_{u\neq0}\frac{Q_C(u)}{\int_Cr^{-2}u^2}=\frac14.
\end{equation}
Their signed divergence satisfies, uniformly near the regular cone,
\begin{equation}\label{eq:cubic-distance-degeneracy}
 q_{X_G}(x)\asymp\frac{\dist(x,C)^3}{r^4}.
\end{equation}
In particular the fields generated by these polynomial potentials do not satisfy a positive
linear distance bound \eqref{eq:linear-distance-divergence}, although they
prove strict minimality and strict stability.
For every sufficiently small \(\eps>0\), the perturbed potential
\(G_\eps=r^{-\eps}G\) has boundary coefficient
\begin{equation}\label{eq:four-degenerate-strictification}
 \beta_{G_\eps}=\eps-\eps^2>0
\end{equation}
and satisfies the distance-weighted comparison of
\cref{prop:distance-weighted-comparison}.
\end{corollary}

\begin{proof}
Along the minimal link the gradient length \(|\nabla G|\) is constant,
since \(h'(0)=1\) in each case and \(A(0)=g^2(1-t^2)\).
Here \(\mu_\ell=0\) in \eqref{eq:mu-window}, so the boundary coefficient
vanishes.  Alternatively this follows from the absence of a linear term in
\(\mathfrak N_\ell[h]\) in the displayed factorizations.
The optimal gap \(1/4\) is \eqref{eq:iso-optimal-gap} at
\(d=4g-2\).
For \((2,3,3)\), equations \eqref{eq:P2} and \eqref{eq:N-z-Q} give
\(\mathfrak N_4[z]=192z^3\).  For the other three potentials,
\eqref{eq:24-certificate}--\eqref{eq:25-certificate} give
\(\mathfrak N_\ell[h]=z^3p(z)\) with \(p>0\) near zero and
\(E(0)>0\).  The spherical function \(z\) is a regular defining
function of the compact minimal link.  Thus \(|z|\asymp\dist(x,C)/r\),
and \eqref{eq:iso-divergence} proves
\eqref{eq:cubic-distance-degeneracy}.
The same factorizations give strict signed divergence everywhere off the
cone.  All hypotheses of \cref{prop:radial-strictification} hold, yielding
\eqref{eq:four-degenerate-strictification} and the claimed distance bound.
\end{proof}

For comparison, Niu \cite[Lemma~5.3]{Niu26} obtains quantitative
subcalibrations from the Hardt--Simon foliation under strict stability
and strict minimality. Here the explicit
potentials suffice, with a radial correction in the four cases above.

\begin{corollary}[Distance-weighted comparison for the nonflat minimizing cases]
\label{cor:iso-distance-weighted}
Every nonflat cone in \cref{thm:iso-polynomial-table} admits a homogeneous
potential, smooth and semialgebraic on \(\R^N\setminus\{0\}\), whose
normalized gradient satisfies \eqref{eq:linear-distance-divergence}.
Except for \eqref{eq:four-degenerate-list}, the polynomial potential in the
table itself has this property.  In the four remaining cases, a sufficiently
small positive rational radial perturbation has it.
\end{corollary}

\begin{proof}
The estimates \eqref{eq:K2-entry} and \eqref{eq:K4-entry} are
strict in their indicated ranges, as are their low-dimensional checks and
endpoint values.  Thus \(Q_\ell>0\) on the entire realizable interval in
all these nondegenerate rows; \eqref{eq:g3-Q} is also strictly positive.
The remaining nondegenerate correction has \(P>0\) by
\cref{prop:17-correction}.
In these cases \(\mathfrak N_\ell[h]=z p(z)\) with \(p>0\) on the
compact interval and \(E>0\).  Near the link \(|z|\) is comparable to
spherical distance, and away from it compactness supplies a positive lower
bound.  Equation \eqref{eq:iso-divergence} and homogeneity prove
\eqref{eq:linear-distance-divergence}.
Apply \cref{cor:four-degenerate-certificates} to the other four cases,
choosing a sufficiently small positive rational \(\eps\). The factor
\(r^{-\eps}=(|x|^2)^{-\eps/2}\), and hence \(G_\eps\), is
semialgebraic on \(\R^N\setminus\{0\}\); see
\cref{prop:radial-strictification}.
\end{proof}

\subsubsection{The two classical exceptions}

Both exceptional cones possess positive Jacobi fields.  Indeed the radial
Jacobi equation is
\begin{equation}\label{eq:iso-radial-Jacobi}
 \gamma(\gamma+d-1)+d(g-1)=0,
\qquad
 \gamma_\pm=-\frac{d-1}{2}\pm\sqrt{D_{g,d}}.
\end{equation}
For \((2,1,5)\) they are \(r^{-2},r^{-3}\); for \((4,1,6)\),
they are \(r^{-6},r^{-7}\). These cases are stable but not minimizing,
completing \cref{thm:classical-iso-classification}.


\begin{corollary}[{Lin's one-sided strict minimality \cite[p.~210]{Lin87}}]
\label{cor:Lin-one-sided}
Let \(x\in\R^2\), \(y\in\R^6\), and set
\[
 E=\{|y|^2<5|x|^2\},\qquad C=\partial\llbracket E\rrbracket,
 \qquad C_1=C\llcorner B_1.
\]
The cone over
\(S^1(\sqrt{1/6})\times S^5(\sqrt{5/6})\) is one-sided strictly
area minimizing in \(\overline E\).
More precisely, it minimizes against finite-perimeter changes
\(H\subset E\) of compact support.  There is \(\Theta>0\) such that
every finite-mass integral current \(T\) with
\[
 \partial T=\partial C_1,\qquad
 \operatorname{spt}T\subset\overline E,\qquad
 \operatorname{spt}T\cap B_\eps=\varnothing
\]
satisfies \(\mathbf M(T)-\mathbf M(C_1)\geq\Theta\eps^7\)
for \(0<\eps<1\).
\end{corollary}

\begin{proof}
Put \(u=|x|^2\), \(v=|y|^2\), \(R=u+v\), and \(F=v-5u\).
Take the radial potential \(G=RF\).  For functions of \(u,v\),
\[
 \Delta P=4uP_{uu}+4vP_{vv}+4P_u+12P_v,\qquad
 \Gamma(P,Q)=4uP_uQ_u+4vP_vQ_v.
\]
Hence
\begin{align*}
 \Delta G&=16(v-8u),\\
 |\nabla G|^2&=16(25u^3+24u^2v+v^3),\\
 \cL(G)&=64R(v-5u)^2(v-17u).
\end{align*}
The last identity follows by inserting the first two in
\[
 \cL(G)=|\nabla G|^2\Delta G-
 \tfrac12\Gamma(G,|\nabla G|^2),
\]
the same radial calculation as
\cref{prop:radial-Q}.
Thus \(X=\nabla G/|\nabla G|\) is smooth off the origin,
equals \(\nu_E\) on the regular cone, and satisfies
\(\diver X<0\) throughout \(E\).
Since \(\cH^{N-1}(\{0\})=0\), the one-sided part of
\cref{thm:paired-subcalibration} gives
\[
 \Per(H;U)-\Per(E;U)
 \geq\int_{E\setminus H}-\diver X\qquad(H\subset E).
\]

Set \(\Theta=\int_{E\cap B_1}-\diver X>0\).
This integral is finite, and homogeneity makes its value over
\(E\cap B_\eps\) equal to \(\Theta\eps^7\).
To pass to currents, use radial projection and an integer-valued BV
filling as in the proof of \cref{thm:strict-minimality-flux}.
The resulting function \(v\) has zero distributional gradient in
the connected open set \(\R^8\setminus\overline E\) and vanishes
there outside a compact set, so it vanishes throughout that set.
It is also constant in \(B_\eps\). Since it vanishes on the
nonempty open set \(B_\eps\setminus\overline E\), this constant
is zero. Hence \(H=\{v\geq1\}\) satisfies \(H\subset E\) and
\(H\cap B_\eps=\varnothing\), up to null sets. Thus \(E\setminus H\)
contains \(E\cap B_\eps\), and the displayed one-sided comparison
proves the claimed mass gap.
\end{proof}
